\chapter{Multiclass networks from first principles}
\label{ch:mnist-mlp}

\begin{learningobjectives}
After this chapter, you should be able to:
\begin{itemize}
  \item formulate handwritten-digit recognition as a controlled multiclass
  experiment;
  \item explain flattening, pixel scaling, and one-hot encoding as explicit
  representation contracts;
  \item derive a dense ReLU network with a logits--softmax--cross-entropy
  output;
  \item motivate Xavier and He initialization from fan-in and signal variance;
  \item derive every term in the vectorized backward pass;
  \item distinguish mini-batch optimization evidence from held-out decision
  evidence; and
  \item diagnose multiclass errors with classwise metrics and individual
  observations.
\end{itemize}
\end{learningobjectives}

\begin{chapterconnection}
Chapter~\ref{ch:neural-networks} derived backpropagation by tracing individual
neurons and then vectorizing them. This chapter puts that mechanism to work on
images. The scientific question is concrete: can a bounded dense network
recognize held-out handwritten digits? The implementation is still small
enough to inspect, but its hidden activations, initialization, multiclass loss,
mini-batches, and evaluation follow modern practice. The next chapter will
replace our manual machinery with a tensor framework without surrendering
these mathematical and evidential contracts.
\end{chapterconnection}

\section{MNIST as experiment, not trophy}

The MNIST database standardized images drawn from NIST handwriting
collections into 60,000 conventional training examples and 10,000 test
examples. Each observation is a $28\times28$ grayscale image with an integer
label in $\{0,\ldots,9\}$. MNIST became historically important because it made
different learning methods easy to compare on a common problem.

Its familiarity is also dangerous. A high benchmark score does not establish
performance on phone photographs, historical manuscripts, a different
writing population, or an operational document pipeline. We use MNIST as a
\emph{mechanism laboratory}: its inputs are rich enough to require learned
representations, while its size and semantics make every stage visible.

The executable laboratory uses a deterministic class-balanced subset of
authentic MNIST containing 2,400 training and 600 held-out images. Versioning
this bounded sample has two engineering advantages: the lesson runs offline,
and continuous integration can execute the same experiment. It has a
statistical cost: classwise estimates have larger sampling uncertainty than
they would on the complete benchmark. Reproducibility and inferential strength
are different properties.

Each digit contributes 240 fitting images and 60 held-out images. Balance
makes overall accuracy easy to read, but it is an experimental design choice,
not a property of every handwriting population. Sixty held-out observations
per class also means one additional class error changes estimated recall by
$1/60\approx0.017$. Apparent differences of one or two percentage points
therefore deserve caution.

\begin{definitionbox}{Experimental unit and prediction target}
The experimental unit is one handwritten digit image. The target is one of ten
mutually exclusive class identities. A model emits a probability distribution
over the fixed class order $0,1,\ldots,9$ and the decision rule selects a class,
usually the largest-probability class.
\end{definitionbox}

Before transformation, inspect image shapes, numeric range, label alignment,
class counts, and representative images. These checks can expose a transposed
axis, corrupt intensities, a missing class, or labels shifted by one position.
A falling loss cannot repair an invalid data contract.

The fitting and held-out partitions have different authorities. Fitting rows
may influence parameters. Validation rows, when present, may influence choices
such as architecture or stopping. Test rows should influence neither; they
estimate performance only after those choices are fixed. The short classroom
experiment displays held-out accuracy during training to teach monitoring, so
it must not be described as a pristine one-time test after the fact.

\conceptcheck{Why is a class-balanced sample convenient for teaching but not a
guarantee that accuracy is the correct metric in a deployed digit-recognition
system?}

\section{From image to feature vector}

For observation $i$, write the image as
\[
  X_i\in\{0,1,\ldots,255\}^{28\times28}.
\]
A dense network consumes a vector, so we flatten the matrix:
\[
  x_i=\operatorname{vec}(X_i)\in\mathbb R^{784},
  \qquad 784=28\cdot28.
\]
With zero-based row and column indices and row-major order, pixel $(r,c)$ maps
to feature index
\[
  j=28r+c.
\]
This formula is part of the interface. A different ordering would be equally
valid during fitting, but inference would have to reproduce it exactly.

The transformation is reversible when its shape and ordering convention are
known. No intensity is numerically discarded. What is discarded is an
\emph{explicit spatial inductive bias}: a dense layer does not know that one
feature corresponds to a pixel immediately above another. It must learn every
useful spatial relationship through independent weights.

\begin{figure}[htbp]
\centering
\begin{tikzpicture}[
  pixel/.style={draw=RiceBlue!65, minimum size=0.25cm, inner sep=0pt},
  feature/.style={circle, draw=RiceBlue, fill=RiceBlue!12, minimum size=0.29cm,
    inner sep=0pt},
  >=Latex
]
  \foreach \r in {0,...,4}{
    \foreach \c in {0,...,4}{
      \pgfmathsetmacro{\shade}{15+mod(17*\r+23*\c,75)}
      \node[pixel, fill=RiceBlue!\shade] at (0.27*\c,-0.27*\r) {};
    }
  }
  \node[above=0.25cm] at (0.54,0) {$28\times28$ image};
  \draw[->, very thick] (1.6,-0.54) -- node[above]{flatten} (3.0,-0.54);
  \foreach \j in {0,...,8}{
    \node[feature] at (3.45+0.36*\j,-0.54) {};
  }
  \node at (6.85,-0.54) {$\cdots$};
  \node[feature] at (7.35,-0.54) {};
  \node[above=0.25cm] at (5.4,0) {$784$ ordered intensities};
\end{tikzpicture}
\caption{Flattening preserves the ordered pixel values but removes the image
grid from the model interface. A convolutional architecture will later encode
locality and weight sharing directly. The small grid is schematic.}
\label{fig:mnist-flattening}
\end{figure}

We scale each intensity by 255:
\[
  \widetilde x_{i,j}=x_{i,j}/255\in[0,1].
\]
This prevents large input magnitudes from overwhelming an initialization
designed under unit-scale assumptions. Unlike standardization by an estimated
mean and variance, division by the known sensor maximum has no fitted state.
It is still part of the inference contract: a deployed caller must apply
exactly the same transformation.

The order of numeric operations matters. Integer images are first converted to
a floating-point dtype and then divided. The resulting feature matrix has
shape $n\times784$ with one observation per row. An executable shape assertion
at this boundary prevents a legal but incorrect reshape from silently changing
which measurements belong to an observation.

\begin{misconceptionbox}
\textbf{Misconception:} flattening ``turns an image into meaningless data.''
\textbf{Correction:} flattening preserves every pixel value and a fixed order.
It removes an architectural statement about two-dimensional locality. A dense
network may learn useful spatial combinations, but it receives no built-in
notion of neighboring pixels.
\end{misconceptionbox}

\section{Class identities and one-hot targets}

The integer label $7$ names a category. It does not imply that class 7 is one
unit farther from class 6 than from class 0. We therefore encode label $k$ as
the $k$th standard basis vector $e_k\in\mathbb R^{10}$:
\[
  7\longmapsto
  (0,0,0,0,0,0,0,1,0,0).
\]
For a batch of $m$ observations, targets form a matrix
$Y\in\{0,1\}^{m\times10}$ whose rows sum to one. Ten columns must be retained
even if one mini-batch happens not to contain every class. Column position is
part of the model's public meaning.

One-hot encoding is useful for deriving the loss. An implementation may accept
integer labels and gather the correct log-probability directly, but it must
preserve the same class-order contract. One-hot targets do not turn one example
into ten examples, nor do they assert ten independent binary tasks.

This distinction determines the output activation. Ten independent sigmoids
would permit probabilities whose sum differs from one and are appropriate for
a multilabel problem in which several labels may be true simultaneously.
Handwritten-digit identity is mutually exclusive, so one softmax couples the
ten scores into one categorical distribution.

\section{A modern bounded architecture}

Let a mini-batch be $A^0=X\in\mathbb R^{m\times784}$. We choose a hidden width
of 64 and compute
\[
  Z^1=A^0W^1+\mathbf1(b^1)^\top,
  \qquad
  A^1=\operatorname{ReLU}(Z^1),
\]
with
\[
  W^1\in\mathbb R^{784\times64},
  \quad b^1\in\mathbb R^{64}.
\]
The output layer computes
\[
  Z^2=A^1W^2+\mathbf1(b^2)^\top,
  \qquad
  W^2\in\mathbb R^{64\times10},
  \quad b^2\in\mathbb R^{10}.
\]
The rows of $Z^2$ are \emph{logits}: unconstrained real-valued class scores.
There is no output ReLU and no independently applied sigmoid. A single softmax
normalizes the mutually exclusive classes:
\[
  p_{i,k}=\frac{\exp z_{i,k}}{\sum_{r=0}^{9}\exp z_{i,r}}.
\]

For one observation, hidden neuron $j$ computes
\[
  z_j^1=\sum_{r=1}^{784}x_rw_{rj}^1+b_j^1,
  \qquad a_j^1=\max(0,z_j^1).
\]
The $j$th column of the batch preactivation matrix collects this same neuron
calculation across observations. Vectorization changes the execution strategy,
not the model. A useful reading discipline is to verify that the scalar
equation, single-observation vector equation, and batch equation describe the
same computation.

\begin{figure}[htbp]
\centering
\begin{tikzpicture}[x=2.6cm,y=0.55cm,>=Latex,
  unit/.style={circle, draw=RiceBlue, fill=RiceBlue!10, minimum size=0.38cm,
    inner sep=0pt}, label/.style={align=center,font=\small\bfseries}]
  \node[label] at (0,4.2) {$784$ scaled\\pixels $A^0$};
  \node[label] at (1,4.2) {$64$ hidden\\features $A^1$};
  \node[label] at (2,4.2) {$10$ logits\\$Z^2$};
  \node[label] at (3,4.2) {$10$ class\\probabilities $P$};
  \foreach \j in {0,...,5}{
    \node[unit] (i\j) at (0,3-\j) {};
    \node[unit] (h\j) at (1,3-\j) {};
  }
  \foreach \j in {0,...,5}{\node[unit] (o\j) at (2,3-\j) {};}
  \foreach \j in {0,...,5}{\node[unit] (p\j) at (3,3-\j) {};}
  \foreach \i in {0,...,5}{
    \foreach \j in {0,...,5}{\draw[RiceBlue!18] (i\i)--(h\j);}
  }
  \foreach \i in {0,...,5}{
    \foreach \j in {0,...,5}{\draw[RiceBlue!18] (h\i)--(o\j);}
  }
  \foreach \j in {0,...,5}{\draw[->,RiceGray,thick] (o\j)--(p\j);}
  \node[below=0.25cm,align=center] at (0,-2) {flatten\\and scale};
  \node[below=0.25cm,align=center] at (1,-2) {dense $+$\\ReLU};
  \node[below=0.25cm,align=center] at (2,-2) {dense\\linear};
  \node[below=0.25cm,align=center] at (3,-2) {softmax};
\end{tikzpicture}
\caption{The from-scratch MNIST classifier. Circles are schematic; ellipses
stand for omitted neurons. The final affine layer emits logits, and softmax
assigns a joint categorical distribution rather than ten unrelated scores.}
\label{fig:mnist-mlp-architecture}
\end{figure}

The model contains
\[
  (784\cdot64+64)+(64\cdot10+10)=50{,}890
\]
trainable scalars. Counting parameters is a useful shape test and a first
measure of capacity, but it does not quantify effective complexity by itself.
The first layer contains more than 98 percent of these parameters because a
dense hidden unit assigns an independent weight to every pixel. This
observation anticipates why spatial weight sharing in a convolution can be
both statistically and computationally useful.

\section{Initialization from variance propagation}

Consider one preactivation
\[
  z_j=\sum_{r=1}^{d}w_{rj}a_r.
\]
Under approximate independence and zero-mean assumptions,
\[
  \operatorname{Var}(z_j)
  \approx d\,\operatorname{Var}(w_{rj})\operatorname{Var}(a_r).
\]
To keep $\operatorname{Var}(z_j)$ on the scale of
$\operatorname{Var}(a_r)$, a linear or approximately symmetric layer suggests
$\operatorname{Var}(w_{rj})\approx1/d$. ReLU maps negative inputs to zero; under
the idealized symmetric-input argument, the second moment of its output is
approximately halved. Compensating gives $2/d$ rather than $1/d$.

With unscaled standard-normal weights, variance grows with fan-in $d$.
Activations can explode, sigmoids can saturate, or numerical ranges can become
unusable before learning begins.

For approximately symmetric signals passed through ReLU, roughly half the
mass is zeroed. He initialization compensates with
\[
  w_{rj}\sim N\!\left(0,\frac{2}{d}\right).
\]
Thus the first layer uses standard deviation $\sqrt{2/784}$. The final linear
logit layer uses Xavier normal scaling,
\[
  w_{rj}\sim N\!\left(0,\frac{1}{d}\right),
\]
giving standard deviation $\sqrt{1/64}$ here. Biases start at zero. Independent
weights already break neuron symmetry, so random biases are not needed.

These calculations are approximations, not distribution-free theorems about
trained networks. Correlations, depth, optimizer dynamics, normalization, and
data geometry matter. Good practice checks initial preactivation scale, the
fraction of active ReLU units, gradient norms, and nonfinite values.

Initialization and reproducibility are related but not identical. A recorded
seed chooses a deterministic pseudo-random stream for a particular
implementation. It does not imply that every numerical library, hardware
accelerator, or parallel reduction will produce bitwise-identical results.
Record the initializer by name and formula in addition to recording the seed.

\begin{warningbox}
Random-normal initialization does not mean ``draw every weight from
$N(0,1)$.'' The variance is part of the algorithm. At fan-in 784, an unscaled
standard normal is about twenty times wider in standard deviation than He
normal initialization.
\end{warningbox}

\section{Softmax cross-entropy as a fused contract}

For one observation with one-hot target $y$, categorical cross-entropy is
\[
  \ell(z,y)=-\sum_{k=0}^{9}y_k\log p_k.
\]
Substituting softmax and using $\sum_k y_k=1$ yields
\[
  \ell(z,y)
  =\log\left(\sum_r e^{z_r}\right)-\sum_k y_kz_k.
\]
This is the log-sum-exp form. Stable code subtracts
$c=\max_r z_r$ before exponentiating:
\[
  \log\left(\sum_r e^{z_r}\right)
  =c+\log\left(\sum_r e^{z_r-c}\right).
\]
Softmax is unchanged because the common factor $e^{-c}$ cancels between
numerator and denominator.

More generally, for any scalar $c$,
\[
  \operatorname{softmax}(z+c\mathbf1)=\operatorname{softmax}(z).
\]
The logits are therefore identifiable only up to a common additive shift.
Probabilities and the largest-probability decision are unchanged. Subtracting
the row maximum exploits this mathematical invariance to obtain numerical
stability; it is not a heuristic alteration of the model.

\begin{proposition}[Softmax--cross-entropy logit gradient]
For $p=\operatorname{softmax}(z)$ and a target distribution $y$ whose entries
sum to one,
\[
  \nabla_z\ell=p-y.
\]
\end{proposition}

\begin{proof}
Differentiate the log-sum-exp expression in coordinate $z_j$:
\[
\begin{aligned}
\frac{\partial\ell}{\partial z_j}
&=\frac{e^{z_j}}{\sum_r e^{z_r}}
  -\frac{\partial}{\partial z_j}\sum_k y_kz_k\\
&=p_j-y_j.
\end{aligned}
\]
Collecting coordinates gives $\nabla_z\ell=p-y$.
\end{proof}

The simplicity of the derivative is not a reason to split softmax and loss in
software. A fused operation states that its inputs are logits, avoids the
logarithm of a rounded zero probability, and can reuse stable intermediates.

\section{Backpropagation for a mini-batch}

Let the objective be the mean of $m$ example losses. The output sensitivity is
\[
  \Delta^2=\frac{P-Y}{m}.
\]
The division by $m$ appears exactly once. If the loss is a mean but the
derivative is treated as a sum, changing batch size silently changes the
effective learning rate. Dividing once in the output sensitivity makes all
subsequent matrix products derivatives of the stated mean loss.

The output parameter derivatives are
\[
  \nabla_{W^2}L=(A^1)^\top\Delta^2,
  \qquad
  \nabla_{b^2}L=(\Delta^2)^\top\mathbf1.
\]
Propagate sensitivity through $W^2$ and apply the ReLU derivative:
\[
  \Delta^1
  =\left(\Delta^2(W^2)^\top\right)
   \odot \mathbf1[Z^1>0].
\]
Then
\[
  \nabla_{W^1}L=(A^0)^\top\Delta^1,
  \qquad
  \nabla_{b^1}L=(\Delta^1)^\top\mathbf1.
\]
Every shape can be audited:
\[
\begin{array}{c|c}
\text{quantity}&\text{shape}\\ \hline
A^0&m\times784\\
Z^1,A^1,\Delta^1&m\times64\\
W^1,\nabla_{W^1}L&784\times64\\
Z^2,P,Y,\Delta^2&m\times10\\
W^2,\nabla_{W^2}L&64\times10.
\end{array}
\]

These shapes explain the transposes. $(A^1)^\top$ has shape $64\times m$ and
$\Delta^2$ has shape $m\times10$, so their product matches $W^2$. Likewise,
$\Delta^2(W^2)^\top$ maps ten output sensitivities back to 64 hidden
sensitivities per observation. A transpose should be justified by this linear
map, not inserted until NumPy accepts the multiplication.

All gradients must be computed from one parameter state before any parameter
is updated. Otherwise an early update changes quantities needed to calculate a
later derivative. Finite-difference checks on a tiny deterministic example
test the analytical implementation independently.

\section{Mini-batch gradient descent}

Full-batch descent uses every training observation for one update. Stochastic
descent uses one. A mini-batch provides a noisy estimate of the full gradient
while exposing matrix operations large enough for efficient numerical kernels.
For batch $B_t$ and learning rate $\eta$,
\[
  \theta_{t+1}
  =\theta_t-\eta\nabla_\theta
   \left(\frac1{|B_t|}\sum_{i\in B_t}\ell_i(\theta_t)\right).
\]

An epoch visits every training observation once after a random permutation.
The random seed, permutation generator, batch size, incomplete final-batch
policy, learning rate, dtype, and initialization are parts of the experimental
record. A seed does not make floating-point computation universally identical
across all hardware, but it removes an important source of uncontrolled
variation.

With $n=2400$ and batch size 64, one epoch contains
$\lceil2400/64\rceil=38$ updates; the final batch contains 32 observations.
Discarding that batch would systematically omit data unless row order were
reshuffled and coverage explicitly managed. Keeping it is simpler, and the
mean-loss gradient makes its update scale comparable to the full batches.

Training loss is \emph{optimization evidence}: it shows whether the procedure
reduces its stated objective on fitting data. Validation or held-out accuracy
is \emph{decision evidence}: it estimates a decision property on unseen data.
Neither alone establishes probability calibration, population transport, or
operational utility.

An epoch-level history is deliberately an aggregate. Batch-level loss is
noisier and can diagnose update instability, while epoch-level loss is easier
to interpret. A production training run should log both only when the added
volume has a stated diagnostic purpose; logging every tensor or observation
can be expensive and can expose sensitive data.

\begin{workedexample}{A bounded handwritten-digit experiment}
The course implementation trains a $784$--$64$--$10$ ReLU network for eight
epochs with batches of 64. Inputs are scaled once, parameters use He and Xavier
initialization under a recorded seed, and a stable fused loss consumes logits.

The experiment first asserts all tensor shapes and verifies that initial
logits are finite. During training it records mean cross-entropy, training
accuracy, and held-out accuracy. After parameters are frozen, it calculates a
ten-by-ten confusion matrix, per-class precision and recall, and a gallery of
high-confidence errors. This sequence answers more than ``what score did the
model get?'' It asks which distinctions remain difficult and whether apparent
confidence agrees with the images a person sees.

The held-out partition is observed during the teaching run to illustrate a
curve. In a genuine model-selection study, that role belongs to a validation
partition; the test set is exposed only after architecture, learning rate, and
stopping policy are frozen.
\end{workedexample}

\section{Confusion is part of the scientific result}

Accuracy is
\[
  \frac{1}{n_{\mathrm{test}}}
  \sum_i\mathbf1[\widehat y_i=y_i].
\]
For class $k$, define
\[
  \operatorname{precision}_k
  =\frac{\mathrm{TP}_k}{\mathrm{TP}_k+\mathrm{FP}_k},
  \qquad
  \operatorname{recall}_k
  =\frac{\mathrm{TP}_k}{\mathrm{TP}_k+\mathrm{FN}_k},
\]
and
\[
  F_{1,k}
  =\frac{2\operatorname{precision}_k\operatorname{recall}_k}
  {\operatorname{precision}_k+\operatorname{recall}_k}.
\]
Precision conditions on the model's class-$k$ decisions; recall conditions on
the true class-$k$ observations. They answer different operational questions
and neither is recovered from overall accuracy alone.

A confusion matrix retains which true digit became which predicted digit.
Rows normalized to sum to one estimate class-conditional recall patterns;
integer counts quantify actual sample evidence. They answer different
questions.

Inspect individual errors after aggregate metrics. A high-confidence mistake
can indicate ambiguous handwriting, an annotation problem, an unusual
translation or stroke, or model overconfidence. Examining only convenient
correct predictions hides failure structure. Examining only the ten worst
examples without aggregate context can exaggerate rare cases. Professional
analysis uses both.

Prediction confidence $\max_k p_k$ is not proof of calibration. Calibration
asks whether events assigned probability near $q$ occur approximately a
fraction $q$ of the time under the relevant distribution. It needs a separate
reliability analysis and enough observations.

The highest-confidence mistakes are a diagnostic sample, not an unbiased
summary of failure frequency. They are chosen precisely because the model was
most certain. Use them to generate hypotheses, then test those hypotheses on a
defined population rather than generalizing from an error gallery.

\begin{professionalbox}
The notebook is an experimental narrative; reusable computations belong in a
tested package. Record data provenance, split identity, code revision,
preprocessing, class order, architecture, initializer, learning rate, batch
size, seed, metrics, and known limitations. A released artifact must reproduce
the same flattening and scaling before applying the learned parameters.

Serialization should package or reference all inference state: parameter
arrays, architecture, dtype, feature ordering, scaling rule, and class order.
A fresh-process test should load the artifact, transform a versioned fixture,
and reproduce a stored probability vector within an explicit tolerance.
\end{professionalbox}

\section{Failure modes and diagnostic tests}

When loss becomes nonfinite, inspect scaled feature range, logits, learning
rate, log-sum-exp stability, parameter norms, and gradients. Accuracy near
$0.1$ suggests a broken ten-class pipeline: verify label alignment, output
width, class order, and whether updates change weights. Falling training loss
with stagnant held-out performance suggests a generalization problem rather
than a derivative failure.

Useful automated tests include:
\begin{itemize}
  \item one-hot rows sum to one and recover their integer labels;
  \item softmax remains finite for large logits and is invariant to a common
  additive shift;
  \item softmax probabilities sum to one;
  \item analytical output and layer gradients match central differences;
  \item initialization standard deviations approximate their fan-in targets;
  \item all parameter and gradient shapes agree; and
  \item a small deterministic multiclass problem can be learned.
\end{itemize}

These tests catch mathematical and interface regressions. They do not establish
that labels measure the intended construct or that a deployment population
matches MNIST.

\section{Exercises}
\begin{enumerate}
  \item Prove that flattening followed by reshaping is invertible when the
  ordering convention and original dimensions are retained. State what the
  dense architecture still fails to encode.
  \item For a $784$--$128$--$64$--$10$ network, write every parameter shape
  and count the trainable scalars.
  \item Derive the softmax Jacobian
  $\partial p_k/\partial z_j=p_k(\mathbf1[k=j]-p_j)$, then recover the fused
  cross-entropy gradient.
  \item Suppose all first-layer weights are sampled from $N(0,1)$. Estimate
  the initial preactivation variance under independent unit-variance inputs
  and compare it with He initialization.
  \item Add a second hidden ReLU layer to the course implementation. Supply
  gradient checks, shape tests, and a controlled comparison that changes only
  depth.
  \item Replace accuracy with a cost matrix in which mistaking 1 for 7 is five
  times as costly as most other errors. Derive the empirical decision cost and
  explain whether retraining is required to change the decision rule.
  \item Reproduce the full study on Fashion-MNIST. Use a validation partition
  for architecture choices, preserve the test set, report per-class evidence,
  and analyze visually plausible confusions such as shirts and pullovers.
  \item Write a pull request adding serialization and a fresh-process inference
  test. The saved contract must include feature shape, class order, scaling,
  architecture, and parameter arrays.
\end{enumerate}

\begin{chapterreview}
An image becomes a dense-network input by an explicit flatten-and-scale
contract. One-hot vectors encode fixed categorical identity; logits are
unconstrained scores; softmax assigns one joint distribution. Pairing logits
with stable cross-entropy gives the output gradient $P-Y$, while vectorized
backpropagation transports that sensitivity through ReLU layers. He and Xavier
random-normal initialization account for fan-in rather than applying arbitrary
noise. Mini-batches make updates practical, but training curves remain
optimization evidence. Held-out confusion patterns and inspected errors turn a
score into a scientific analysis. The manual implementation makes each of
these decisions visible before a framework automates their execution.
\end{chapterreview}

\begin{notebooklink}
Lecture 12 notebook 02 performs the complete authentic-MNIST experiment with a
tested NumPy implementation. It visualizes flattening, audits initialization,
derives and checks the fused loss, trains with mini-batches, studies confusion
patterns, compares hidden widths, and leaves Fashion-MNIST as an independent
applied investigation.
\end{notebooklink}
